\section{Practical Implications: Bias as a Controllable Source of Instability}
\label{sec:implications}

The preceding sections established three key empirical and theoretical facts:

\begin{enumerate}
    \item Activation bias is statistically significant and structurally coherent,
    while gradient output mean remains at noise level.
    \item Extreme activation outliers are overwhelmingly dominated by the bias
    component rather than spike variance or tail variability.
    \item Bias inflates both blockwise quantization scale and adaptive optimizer
    second-moment statistics.
\end{enumerate}

These observations allow us to reinterpret several long-standing training
pathologies through a unified lens: many instabilities attributed to spectral
anisotropy are in fact primarily driven by activation bias.

We now discuss the implications for quantization and optimization.

\subsection{Quantization: Bias as Dynamic-Range Inflation}

Low-bit quantization operates under a fundamental constraint: each block of
features shares a common scaling factor. The effective resolution allocated to
each element is therefore inversely proportional to the blockwise scale.

From Appendix~\ref{app:quant_bias}, if a block contains coordinates with large
$|\mu_j|$, the scale concentrates around
\[
\alpha(x;B) \approx \max_{j \in B} |\mu_j|,
\]
rather than the intrinsic stochastic variability level
$O(\sigma \sqrt{\log b})$.

This inflation has two consequences:

\begin{itemize}
    \item The quantization step size increases proportionally.
    \item The long-tail semantic component $s_j$ is compressed relative to
    its intrinsic scale.
\end{itemize}

Empirically, mean removal reduces blockwise $p99.9$ magnitudes by
approximately a factor of two in early layers, consistent with the
high-probability bounds derived in Appendix~\ref{app:quant_bias}.
The resulting reduction in quantization step size directly improves effective
tail signal-to-noise ratio (SNR).

Importantly, this mechanism does not require full spike whitening.
Bias removal alone addresses the dominant source of dynamic-range inflation.
Spike whitening may provide further improvements when spike variance remains
substantial, but it is not necessary to eliminate the principal outlier driver.

\subsection{Optimization: Bias as Second-Moment Inflation}

Adaptive optimizers such as Adam maintain elementwise second-moment estimates
\[
v_{ab} = \mathbb{E}[g_{ab}^2],
\qquad
g_{ab} = \delta_a x_b.
\]

Appendix~\ref{app:gradient_proof} shows that
\[
v_{ab}
=
\mu_b^2 \mathbb{E}[\delta_a^2]
+
\mathbb{E}[\delta_a^2 s_b^2],
\]
under mild independence assumptions.

When $\mu_b^2$ dominates $\mathbb{E}[s_b^2]$, the second moment becomes
bias-driven. This inflates $v_{ab}$ for specific coordinates, thereby
suppressing their effective learning rates via the update
\[
\Delta W_{ab} \propto \frac{g_{ab}}{\sqrt{v_{ab}}}.
\]

Empirically, mean removal reduces the extreme tail of $v$ by factors ranging
from $4\times$ to nearly $50\times$ in early layers, while leaving median
values largely unchanged. This confirms that bias primarily affects a small
subset of coordinates that nonetheless dominate adaptive learning-rate
suppression.

Unlike gradient-output mean (which is statistically noise-level), activation
bias propagates multiplicatively through the outer-product structure
$g_i = \delta_i x_i^\top$ and therefore directly controls second-moment
inflation.

\subsection{When Is Mean Removal Sufficient?}

The experiments suggest a layered interpretation:

\begin{itemize}
    \item \textbf{Early layers / early training:} bias dominates both
    outliers and second-moment inflation.
    Mean removal alone resolves most instability.
    
    \item \textbf{Later layers / late training:} spike variance persists even
    after bias removal.
    Additional spectral whitening may further flatten the spectrum.
\end{itemize}

Thus, bias removal can be viewed as the minimal structural correction.
Spike whitening becomes an incremental refinement rather than the primary fix.

\subsection{Relation to Existing Spectral Methods}

Full SVD-based spike removal eliminates both bias and spike variance.
Our analysis clarifies that a substantial portion of the benefit previously
attributed to spike suppression may arise from removing the coherent mean
component.

This observation has two practical consequences:

\begin{enumerate}
    \item In low-bit training, mean removal can replace expensive SVD
    decompositions when bias dominates dynamic-range inflation.
    
    \item In optimizer design, centering activations may mitigate second-moment
    inflation without requiring explicit spike tracking.
\end{enumerate}

These implications motivate a re-examination of spectral regularization
techniques: in many regimes, controlling bias may be sufficient to stabilize
training.

\subsection{Summary}

The empirical and theoretical analyses converge on a consistent picture:

\begin{itemize}
    \item Bias is the primary driver of extreme activation outliers.
    \item Bias propagates multiplicatively into gradient second moments.
    \item Removing bias reduces both quantization instability and optimizer
    suppression.
    \item Spike variance remains but plays a secondary role relative to bias.
\end{itemize}

These results reposition mean removal from a heuristic preprocessing step to a
structurally justified intervention grounded in activation geometry.